\section{新补入的实验和观测：Hall特征也有模型歧义}
\subsection{MAIZE导引场实验：四极密度并非直接测得四极磁场}
Varnish等\citep{F703}在约500 kA、200 ns的MAIZE电流驱动双碳丝阵列中，以阵列倾角调节$B_g/B_{\rm rec}=0,0.4,1$。两正交视线的1064 nm、2 ns干涉诊断显示，导引场下密度层产生旋转和四极密度特征；不同时间的重复发次提供了其持续性的证据。典型$d_i$估计约0.5 mm，装置约$40d_i$，不是尺寸无限大的理想平板。

该文的三维计算是GORGON\emph{电阻MHD}，100 $\mu$m网格、$480\times400\times280$，不是自洽Hall计算。扭曲层几何可解释部分投影，却不能复现四极密度；作者进一步给出预设密度的Hall启发式模型，能得到更相似投影。这使Hall解释成为值得验证的候选，但没有直接测量层内Hall磁场、完整流速和温度来唯一确认它。两条视线也不足以唯一反演三维密度；模拟密度比实验高约2--3倍，零点校正依赖质量烧蚀对倾角不敏感等假设。

因此应写“与Hall机制相容的四极密度结构”，而不是“实验已经自洽验证了三维Hall重联模型”。密度团块也不自动等于闭合磁通岛。最直接的新研究是使用同一驱动、辐射/电离模型及边界加入Hall项，再合成相同干涉视线，检查新增四极信号与实测幅度是否同时匹配。

\subsection{MMS冷离子：同一位置的热、冷离子不一定同时退磁}
André等\citep{L227}分析2015年11月12日MMS磁层顶分离线事件。四星间距约10 km，电子/离子采样30/150 ms；跨层宽度估计70--100 km，冷离子约50 eV、回旋半径约25 km，而较热离子的回旋尺度约100--320 km。冷离子仍磁化、热离子已退磁，因此使用单一离子温度会丢失不同尺度响应。

法向电场约10 mV/m，压力项约2--3 mV/m；冷离子密度分数约一半，其对流项能与Hall项相当。这里测的是分离线法向场，不能改称X线重联电场。电场处理中还用了$\E\cdot\B=0$约束和约2 mV/m偏置校正，压力/电流梯度来自四星平均；这些诊断条件应与结论一起保留。此单事件支持多离子模型的必要性，没有给出冷离子对全局率的普适修正因子。

\subsection{MRX部分电离：Hall形态保留，喷流却明显减速}
Lawrence等的MRX实验\citep{L190}保持离子密度和电子温度大致相近，提高氦中性密度。Hall四极场相对上游场的形状和幅度接近不变，离子喷流速度却降低并接近按总质量密度定义的Alfvén速度。喷流通道宽度没有按无限域强耦合理论预期扩大，导致以离子Alfvén速度归一的重联率下降；换用总密度速度则近常数。

这不是简单的“电阻增大所以速率减小”：较高密度区电阻率平台或上升时，率仍下降，离子动量与中性拖曳是重要限制。最高中性密度时预期有效Hall尺度约83 cm，超过实际设备尺度，边界可能阻止通道扩大。中性密度约有50\%误差，中性压力未直接测量，所以这一边界解释仍需联合模拟验证。

\subsection{MST比较：应按松弛前态分类，再谈模型是否吻合}
扩展MHD研究\citep{F798}比较了MST电流与平行速度在松弛期间的改变。首次准单螺旋态崩解与后续多螺旋松弛给出不同符号关系；后者才更接近常规实验集合平均。该比较是模拟作者对既有实验数据的分析，不是本项目另行精读了其全部实验引用。也不能把模拟壁面的动量注入解释直接搬到真实装置，原文对此明确保留。
